\section{Palindrome Queries}
\label{sec:cses-palindrome-queries}

\problemstatement{Given a string, process two kinds of queries: (1) change
the character at position $k$; (2) ask whether the substring $[a,b]$ is a
palindrome. Answer all type-2 queries online.}

\begin{patternbox}
\emph{``Update a character, then test a substring for palindrome''}
combines point updates with range equality checks -- exactly what
\textbf{polynomial hashing over a Fenwick tree} supports. Maintain a
forward hash and a reversed hash; the substring is a palindrome iff the
two agree.
\end{patternbox}

\subsection*{Two hashes, both point-updatable}

Store the string's characters weighted by powers of a base $B$ in two
Fenwick (BIT) trees: one for the forward direction, one for the reversed
direction. The hash of $[a,b]$ read forwards is a normalised range sum
from the forward BIT; the hash of $[a,b]$ read backwards is the
corresponding range sum from the reversed BIT. Normalising both to the
same starting power lets them be compared directly. A character update is
two point updates (one per tree). The substring is a palindrome exactly
when the forward and reversed hashes match.

\begin{ideabox}[Palindrome = Forward Hash Equals Reversed Hash]
A substring reads the same both ways iff its forward and backward hashes
coincide. Encoding each character as $c\cdot B^{\text{pos}}$ in a BIT makes
both hashes range sums, so point updates and range queries are each
$O(\log n)$. Use two independent moduli/bases to make hash collisions
negligible.
\end{ideabox}

\subsection*{Algorithm}

\begin{enumerate}
  \item Precompute powers of $B$; insert each character into a forward
        and a reversed Fenwick tree.
  \item Update: point-update both trees at the changed position.
  \item Query $[a,b]$: fetch the forward and reversed range hashes,
        normalise to a common power, compare; equal $\Rightarrow$
        palindrome.
\end{enumerate}

\subsection*{Dry run}

\dryrun{$s=$``aybabtu''; query palindrome on ``bab'' ($[2,4]$).}

\begin{itemize}
  \item Forward hash of ``bab'' and reversed hash of ``bab'' both encode
        $b,a,b$ in mirror order and match.
  \item After updating position $3$ (`a'$\to$`x'), ``bxb'' still matches
        (palindrome); ``bxa'' would not.
\end{itemize}

\subsection*{C++ implementation}

\begin{lstlisting}
#include <bits/stdc++.h>
using namespace std;
typedef unsigned long long ull;
const ull B = 131;

struct BIT {
    int n; vector<ull> tree;
    BIT(int n) : n(n), tree(n + 1, 0) {}
    void update(int i, ull val) { for (++i; i <= n; i += i & -i) tree[i] += val; }
    ull query(int i) { ull s = 0; for (++i; i > 0; i -= i & -i) s += tree[i]; return s; }
    ull range(int l, int r) { return query(r) - (l ? query(l - 1) : 0); }
};

int main() {
    int n, q;
    string s;
    cin >> n >> q >> s;
    vector<ull> pw(n + 1, 1);
    for (int i = 1; i <= n; i++) pw[i] = pw[i - 1] * B;

    BIT fwd(n), rev(n);
    for (int i = 0; i < n; i++) {
        fwd.update(i, (ull)s[i] * pw[i]);          // weight by position
        rev.update(i, (ull)s[i] * pw[n - 1 - i]);
    }

    while (q--) {
        int type; cin >> type;
        if (type == 1) {
            int k; char c; cin >> k >> c; k--;
            fwd.update(k, ((ull)c - s[k]) * pw[k]);
            rev.update(k, ((ull)c - s[k]) * pw[n - 1 - k]);
            s[k] = c;
        } else {
            int a, b; cin >> a >> b; a--; b--;
            ull h1 = fwd.range(a, b) * pw[n - 1 - b];       // normalise
            ull h2 = rev.range(a, b) * pw[a];
            cout << (h1 == h2 ? "YES" : "NO") << "\n";
        }
    }
    return 0;
}
\end{lstlisting}

\begin{complexitybox}
\textbf{Time Complexity.} $O((n+q)\log n)$ -- each update and query is
$O(\log n)$. \textbf{Space Complexity.} $O(n)$ for the two Fenwick trees
and power table.
\end{complexitybox}

\begin{takeawaybox}
Dynamic substring equality is polynomial hashing plus a point-updatable
Fenwick tree; a palindrome test is just ``forward hash equals reversed
hash.'' For safety against adversarial collisions, run two independent
hashes. This forward/reversed-hash pattern answers many online substring
comparison queries.
\end{takeawaybox}
